% Part: normal-modal-logic
% Chapter: syntax-and-semantics
% Section: normal-modal-logics

\documentclass[../../../include/open-logic-section]{subfiles}

\begin{document}

\olfileid{nml}{syn}{nor}

\olsection{正规模态逻辑}

模态逻辑可被视为任何具有良好性质的公式集合。某些模态逻辑尤其受到关注，既因为它们有着有趣的解释与应用，但主要是因为它们在理论上已被透彻理解。这些就是所谓的正规模态逻辑。

\begin{defn}\ollabel{defn:modal-logic}
一个\emph{模态逻辑}~$\Log L$ 是基本模态语言的公式集合，它满足：
\begin{enumerate}
\item 包含所有命题重言式；
\item 在统一代入下封闭；且
\item 在分离规则下封闭，即，只要 $!A$ 和 $(!A \lif !B) \in \Log L$，则 $!B \in \Log L$。
\end{enumerate}
\end{defn}

\begin{ex}
满足此定义因而可算作模态逻辑的公式集合（尽管并不十分有趣）包括：
\begin{enumerate}
\item 所有重言式的集合
\item 所有公式的集合（不一致的模态逻辑）
\end{enumerate}
\end{ex}

\begin{defn}\label{defn:normal-modal-logic}
一个模态逻辑~$\Log L$ 是\emph{正规的}，当且仅当它同时满足：
\begin{enumerate}
\item 包含
\begin{align}
\tag{K} \Box(p \lif q) & \lif (\Box p \lif \Box q) \\
\tag{Dual} \Diamond p & \liff \lnot \Box \lnot p
\end{align}
\item 在必然化规则下封闭，即，只要 $!A \in \Log L$，则 $\Box !A \in \Log L$。
\end{enumerate}
\end{defn}

我们将看到定义模态逻辑的两种主要方式。一种是证明论的方式，即将其定义为在特定推导系统中可推导出的公式集合。另一种是模型论的方式，即将其定义为在特定模型类中有效的公式集合。这与普通命题逻辑（乃至一阶逻辑）中的情形相对应。在那里，我们同样可以用两种不同的方式刻画一个公式集合，例如，既可以将其定义为在所有真值赋值下均为真的重言式集合，也可以将其定义为某个推导系统中所有可推导公式的集合。对于正规模态逻辑，利用关系模型和公理化（希尔伯特型）证明系统，这种以双重方式刻画模态逻辑的做法是可行的。
\end{document}